\section{Harmonic collapse and complete rescaled limits}\label{sec:harmonic-limits}

The classes calculated above determine exact periods.
To determine diameters, one must also compare paths
in the full compact moduli space. We first construct
compatible local comparisons and explain why short
ambient paths remain in them. This permits a relative
estimate even when one exceptional plane is much
smaller than the other.

\subsection{Exact periods of the reconstructed cycles}
Write
\[
 b_\eta=\left(0,[\omega_\eta],
       \frac{c_1(v)\cdot[\omega_\eta]}{\rank v}\right).
\]
Equation~\eqref{eq:harmonic-class} gives
$[\Omega_\eta]=2\pi^2\theta_v(b_\eta)$.
In the interacting family
$(a_i,b_\eta)=\deg_\eta L_i-\mu_\eta(v)
=-\zeta_i/(2\pi)$. Pairing with
\eqref{eq:interacting-curve-classes} proves
\begin{equation}\label{eq:exact-cycle-volumes}
 \begin{gathered}
 \operatorname{Area}(c_0)=\pi a,\qquad
 \operatorname{Area}(c_2)=\pi b,\qquad
 \operatorname{Area}(f)=\pi(a+b),\\
 \Vol(P_0)=\frac{(\pi a)^2}{2},\qquad
 \Vol(P_2)=\frac{(\pi b)^2}{2}.
 \end{gathered}
\end{equation}
The volume formula is $\int_P\Omega_\eta^2/2$.
These are exact identities for the actual harmonic
metric; they do not require its induced plane
metric to be Fubini--Study. The reduced
specialization of curves is reflected in exact
additivity of their areas.

For the isolated example put
$\kappa_t=44\pi|t|/3$. Its signed line class
\eqref{eq:isolated-classes} similarly gives
\begin{equation}\label{eq:isolated-periods}
 \operatorname{Area}(\ell_\pm)=\pi\kappa_t,
 \qquad \Vol(P_\pm)=\frac{(\pi\kappa_t)^2}{2}.
\end{equation}
The same $\kappa_t$ is the absolute moment level
in \eqref{eq:isolated-level}.

\subsection{Compatible comparisons over the entire exceptional set}
At the deeper point use the fixed background
$E_0=L_0\oplus L_1\oplus L_2$. Its chart
$\Phi_\eta$ has the full model
\eqref{eq:three-vertex-model} and the estimate
\eqref{eq:actual-quotient-comparison}, with radius
$r=\|x\|$. The common radius and constant are
independent of ratios of positive eigenvalues
of the orbit operator $C_x^*C_x$.
Section~\ref{sec:compact-interaction} proves that
the chart covers whole fibres and commutes with
the contraction to $Z$.

For $F$ in a compact subset of $S^\circ$, local
universal families give smoothly varying central
HE connections on $L_1\oplus F$. The elliptic
complement has a uniformly positive lower bound
on that compact set. Harmonic projections and
the finite-regularity estimates vary smoothly,
so finitely many charts have uniform constants.
Their full harmonic representations are
\begin{equation}\label{eq:moving-normal-model}
 U\oplus\bigl(\C^2\oplus(\C^2)^*\bigr),
 \qquad \dim_\C U=2,
 \qquad \kappa=2\pi(2\mu_\eta(v)-\deg_\eta F)=-\zeta_1.
\end{equation}
The effective group is a circle, trivial on $U$.
The factor two in the target comes from the rank
of $F$; the multiplicity-space moment still uses
ordinary trace.

Near $o$ these two-factor gaps need not remain
positive. Instead choose one fixed smaller
neighborhood of $o$ covered by the three-factor
chart. Its coverage includes moving fibres whose
base tends to $o$, by the normalized-map
compactness argument of
Section~\ref{sec:compact-interaction}. The
remaining part of the contracted locus is compact
in $S^\circ$ and admits the preceding finite cover.

On an overlap the transition is exactly
$h_{ij,\eta}=\Phi_{j,\eta}^{-1}\Phi_{i,\eta}$.
Both charts classify the same holomorphic
bundles; HE and finite balancing change their
presentations, not their isomorphism classes.
The transitions are holomorphic on the complex
quotient charts and commute with the contraction.
A base-line twist of a local universal family
does not change either assertion.
For $\varepsilon_i=C_i(|\eta|+r_i^2)<1/2$,
comparison with the same actual metric gives
\begin{equation}\label{eq:overlap-metric}
 \frac{1-\varepsilon_i}{1+\varepsilon_j}g_i
 \le h_{ij,\eta}^*g_j
 \le\frac{1+\varepsilon_i}{1-\varepsilon_j}g_i.
\end{equation}
This relative statement remains valid when
positive orbit eigenvalues become arbitrarily
small. A bound for ordinary derivatives of
balancing gauges is not required.

\subsection{From tensor comparison to ambient distances}
Equip each $\Q_\zeta$ with the length distance
$d_\zeta$ obtained from continuous paths which
are smooth in representation-type strata on a
finite subdivision, with finite one-sided
lengths at endpoints. The complex zero fibre
has the module-simple point
\eqref{eq:simple-point}. The finite intrinsic
metric theorem \cite[Section~5]{wuebben14}
therefore applies: these are proper geodesic
spaces with the compact quotient topology,
their distances vary continuously with balanced
representatives and levels, and every target
point is approximated at nearby levels.
The quotient norm radius is 1-Lipschitz.

Shrink the actual chart to have squared-metric
bounds $1/2$ and 2 on $r<2R$. Closed radial
sublevels are compact in the quotient and their
images are closed in the compact Hausdorff
moduli space. An actual path starting at radius
$r_0<R$ and leaving the radius-$R$ image must
pay at least
\begin{equation}\label{eq:first-exit}
 \frac{R-r_0}{\sqrt2}
\end{equation}
before its first exit. A segment of length $L$
remaining in the chart has radius at most
$r_0+\sqrt2L$. This applies to paths in the entire
moduli space and excludes a shortcut through
a distant part of that space.

Here is the resulting complete rescaling
statement, with its basepoint data made explicit.
Let $\eta_n\to0$ be generic, $\delta_n\downarrow0$,
and
\begin{equation}\label{eq:first-rescaling-data}
 \nu_n=\frac{\zeta(\eta_n)}{\delta_n^2}\longrightarrow\nu.
\end{equation}
Let $p_n$ be balanced representatives at levels
$\nu_n$ converging, after unitary adjustment,
to a representative $p$ at level $\nu$. Then
\begin{equation}\label{eq:first-rescaling-limit}
 \bigl(M_{\sigma(\eta_n)},\delta_n^{-1}d_{g_{\eta_n}},
       \Phi_{\eta_n}([\delta_np_n])\bigr)
 \longrightarrow(\Q_\nu,d_\nu,[p])
\end{equation}
in pointed Gromov--Hausdorff distance.

For completeness, homogeneity identifies the
current finite model with $\delta_n\Q_{\nu_n}$.
Finite near-minimizers on a bounded rescaled
ball have bounded radius, by its 1-Lipschitz
property, and scale into the actual chart.
Conversely an actual competitor of length
$O(\delta_n)$ cannot pay the fixed exit cost
\eqref{eq:first-exit}; its radius is
$O(\delta_n)$. Applying
\eqref{eq:actual-quotient-comparison} to both
types of competitors and taking both infima
gives relative distance error
\begin{equation}\label{eq:rescaled-relative-error}
 O_L(|\eta_n|+\delta_n^2)
\end{equation}
on each bounded rescaled ball. This comparison
also holds for pairs much closer than
$\delta_n$. The inverse chart covers every
actual bounded ball, and the model distance
continuity identifies the full complete target.
The estimate compares the actual metric with
the \emph{current} quotient $\Q_{\nu_n}$;
it does not prescribe a rate for
$\Q_{\nu_n}\to\Q_\nu$.

At a fixed point of $S^\circ$ the same proof
uses \eqref{eq:moving-normal-model}. It gives
the complete product of the flat harmonic
tangent space $U$ and the Eguchi--Hanson
quotient $T^*\PP^1$ at nonzero normalized level,
or $U\times(\C^2/\{\pm1\})$ at zero level.
The statements are uniform on compact subsets
of $S^\circ$.

\subsection{Two scales in the deepest fibre}
Let $C_i(a,b)$ denote the finite planes identified
with $P_i$, and put $\rho=a+b$. Their balanced
radii and finite ambient diameters are
$O(\sqrt\rho)$, by the finite core bound
\cite[Section~6]{wuebben14}. The preceding
first-exit argument, now at scale $\sqrt\rho$,
gives
\begin{equation}\label{eq:relative-plane-diameter}
 \diam_{g_\eta}P_i
  =\sqrt\rho\,(1+O(\rho))
        D_i(a/\rho,b/\rho),
\end{equation}
where $D_i(A,B)$ is the ambient diameter of
the indicated plane in the complete quotient
$\Q_{(A,-A-B,B)}$. The bounded inverse of the
linear parameter map makes $|\eta|=O(\rho)$
in this chamber. The error is relative even
when $D_i\to0$: shorter competitors stay in
the same radius-$O(\sqrt\rho)$ region, so no
division by a weak diameter occurs.

An explicit calculation separates the scales.
On $C_0$, use \eqref{eq:plane-zero-section} and
fix $w$ positive real. Before taking the
remaining common phase quotient of $(z,\varphi)$,
minimizing over the eliminated phase gives
\begin{equation}\label{eq:plane-metric-form}
 |dz|^2+|d\varphi|^2
 +\frac{(\operatorname{Re}\langle\varphi,d\varphi\rangle)^2}
             {b+|\varphi|^2}
 -\frac{(\operatorname{Im}\langle\varphi,d\varphi\rangle)^2}
             {b+2|\varphi|^2}.
\end{equation}
It lies between one half and twice the flat
form on $(z,\varphi)$. Infimizing over the
remaining phase preserves these inequalities.
Thus the induced plane metric is bounded
between one half and twice Fubini--Study at
line area $\pi a$. Its intrinsic diameter is
at most $\pi\sqrt a/\sqrt2$, an ambient upper
bound as well.

For a lower bound use the points with $z=0$,
$\varphi=\sqrt a(1,0)$ and $\sqrt a(0,1)$,
and $w=\sqrt{a+b}$. Any pair of their compact
group translates has Euclidean separation at
least $\sqrt{2a}$. Quotient length distance
dominates Euclidean orbit separation. Hence
\begin{equation}\label{eq:uniform-plane-bounds}
 \sqrt{2a}\le\diam_{d_\zeta}C_0(a,b)
       \le\frac\pi{\sqrt2}\sqrt a.
\end{equation}
Interchanging vertices gives the corresponding
bounds with $b$ for $C_2$. Combined with
\eqref{eq:relative-plane-diameter}, this proves
the ratio-independent bounds of
Theorem~\ref{thm:main}.

Suppose $a/b\to0$ and $\rho\to0$. After dividing
by $\sqrt b$, put $\epsilon=a/b$. The entire
plane satisfies
$|z|^2+|\varphi|^2=\epsilon$ and
$w=\sqrt{1+|\varphi|^2}=1+O(\epsilon)$.
It approaches the wall point $w=1$ uniformly.
Its stabilizer separates vertex~0 from the
spherical pair $(1,2)$. The full normal
representation has positive coordinates
$(z,\varphi_1,\varphi_2)$, their negative
duals, and effective level $\epsilon$.
There is no tangent factor. A second rescaling
by $\sqrt\epsilon$ gives the unit-level
Calabi quotient on $T^*\PP^2$, and the whole
plane converges to its zero section.

In that complete quotient
\[
 \{(q,p):p(q)=0,\ |q|^2-|p|^2=1\}/\mathrm U(1),
\]
the map $[q,p]\mapsto[q]$ is a retraction to
the zero section. Since $|q|\ge1$, its
Fubini--Study differential is bounded by the
ambient flat norm. Infimizing over vertical
directions makes it 1-Lipschitz on the
quotient. Ambient and intrinsic zero-section
distances therefore agree, and their diameter
is $\pi/2$. The relative comparisons at both
scales prove
\begin{equation}\label{eq:separated-diameter}
 \frac{\diam_{g_\eta}P_0}{\sqrt a}\longrightarrow\frac\pi2
       \quad(a/b\to0),
 \qquad
 \frac{\diam_{g_\eta}P_2}{\sqrt b}\longrightarrow\frac\pi2
       \quad(b/a\to0).
\end{equation}

\subsection{Successive complete limits}
The two-scale calculation is a particular case
of a statement with moving centers. The necessary
data include both the residual stability level
and the normal basepoint.

\begin{theorem}\label{thm:successive-limits}
Start with the actual rescaling
\eqref{eq:first-rescaling-data}--\eqref{eq:first-rescaling-limit}.
Let the current finite levels $\nu_n$ converge
to $\nu_*$ and the current basepoints converge
to $[x_*]\in\Q_{\nu_*}$. In the full normal
slice at $x_*$ write the representation as
$U\oplus W'$ with compact stabilizer $H$,
where $H$ acts trivially on $U$. Write the
balanced coordinates of the basepoints as
$(u_n,z_n)\to0$. Suppose $\sigma_n\downarrow0$ and
\begin{equation}\label{eq:successive-data}
 \frac{(\nu_n-\nu_*)|_{\Lie H}}{\sigma_n^2}
       \longrightarrow\lambda,
 \qquad \frac{z_n}{\sigma_n}\longrightarrow z_*
\end{equation}
after compact group adjustment. At physical
scale $\delta_n\sigma_n$ the complete pointed
limit is
\begin{equation}\label{eq:successive-product}
 U\times(W'///_\lambda H),\qquad (0,[z_*]).
\end{equation}
Here $///$ denotes the hyperk\"ahler quotient
with complex moment zero and real moment
$\lambda$. The statement may be iterated
through any fixed finite succession satisfying
\eqref{eq:successive-data} at each stage.
The physical scale is the product of the
successive scales. No bound on $|u_n|/\sigma_n$
is required.
\end{theorem}

\begin{proof}
The full normal-slice and moving-center
theorems \cite[Sections~4 and~6]{wuebben14}
give the comparison after the change
$(u,z)\mapsto(u_n+\sigma_nu,\sigma_nz)$.
On bounded normal balls its relative
squared-metric error is
\begin{equation}\label{eq:successive-error}
 O_{x_*}(|\nu_n-\nu_*|+|u_n|+\sigma_n R)=o(1).
\end{equation}
Translation removes the tangential motion
and does not change a moment level. The
outer error \eqref{eq:rescaled-relative-error}
is also multiplicative, so it adds to this
error and is not amplified by dividing by
$\sigma_n$.

Choose a fixed normal sublevel whose twice-sized
neighborhood lies in the slice. Its exit cost,
divided by $\sigma_n$, tends to infinity.
An actual path of length $O(\delta_n\sigma_n)$
therefore remains in both the outer chart and
the inner slice; its translated normal radius
is $O(\sigma_n)$. Conversely every bounded
normal-model near-minimizer maps into these
charts. Taking both length infima proves
relative distance comparison, and the common
inverse gives full ball coverage.

The existence of a module-simple point passes
to every full normal Ext representation
\cite[Section~6]{wuebben14}. Its finite metric
theorem applies to the full $U\oplus W'$;
the product decomposition then gives the
complete target \eqref{eq:successive-product}.
For finitely many stages the relative errors
sum to $o(1)$, and the successive exit estimates
and inverse charts give the asserted iteration.
Each limiting center has its own fixed radius
and constants. The statement requires the
convergent data at every stage and makes no
assertion about an infinite succession.
\end{proof}

For this quiver the nonzero $\zeta_0$ and
$\zeta_2$ walls have the isolated
$T^*\PP^2$ normal model. On a nonzero
$\zeta_1$ wall the tangent factor is $\C^2$
and the transverse model is $T^*\PP^1$.
If the internal scale in that pair-moduli
surface also vanishes, the same full normal
iteration applies. This describes approaches
from all six chambers without a lower bound
on a ratio of positive scales.

\subsection{The isolated and moving fibre constants}
In the isolated example the model plane is
the zero section of $T^*\PP^2$ at level
$\kappa_t$. Its induced metric is exactly
$\kappa_t g_{\rm FS}$, with line area $\pi\kappa_t$.
The actual comparison has error $O(|t|)$
on that entire section. The first-exit
argument and the 1-Lipschitz retraction give
\begin{equation}\label{eq:isolated-diameter}
 \diam_{g_t}P_\pm=\frac\pi2\sqrt{\kappa_t}
                       (1+O(|t|)).
\end{equation}
Dividing distances by $\sqrt{\kappa_t}$ gives
the complete unit Calabi quotient, with the
zero section as the limit of the actual plane.

For a moving fibre $F_y$ over $y\in S^\circ$,
the same argument in the two-arrow model
gives
\begin{equation}\label{eq:moving-diameter}
 \operatorname{Area}(F_y)=\pi|\zeta_1|,
 \qquad
 \frac{\diam_{g_\eta}F_y}{\sqrt{|\zeta_1|}}
       \longrightarrow\frac\pi2.
\end{equation}
The limit is through nonzero transverse
levels as $\eta\to0$, uniformly for $y$
in a compact subset of $S^\circ$.
The complete normal limit is Eguchi--Hanson
times the harmonic tangent plane. When
$y\to o$, the fixed three-factor chart and
the retained rescaled basepoint supply the
applicable limit instead. These are the
same fibres that specialize in
\eqref{eq:proper-specialization}.
